% Questions de cours issues de 8-ModelisationSA2023.docx.
\begin{enumerate}
\item indiquer l'intérêt d'utiliser la transformée de Laplace ;
\item à quel rapport correspond la fonction de transfert d'un système ou d'un constituant~? Donner sa forme générale en précisant gain statique, classe et ordre~;
\item donner la méthode permettant d'obtenir cette forme pour une fonction de transfert~;
\item donner la condition fondamentale de stabilité ;
\item quel type de pôle génère des oscillations~;
\item donner le théorème de la valeur finale, et celui de la valeur initiale ;
\item donner les transformées de Laplace des entrées tests ;
\item donner la valeur de l'erreur statique absolue et relative d'un modèle stable non perturbé soumis à un échelon~;
\item donner la valeur finale d'un modèle stable non perturbé soumis à un échelon~;
\item qu\textquotesingle est-ce qu'un système asservi ?
\item représenter la structure générale d'un schéma-bloc d'une grandeur asservie ;
\item expliquer chaîne directe et chaîne de retour ;
\item donner la relation entre les fonctions de transfert de l\textquotesingle IHM et du capteur ;
\item expliquer la simplification de blocs en série, en parallèle, en boucle fermée.
\item représenter la structure générale d'un schéma-bloc d'une grandeur asservie perturbée ;
\item par quoi se traduit une équation à 3 variables dans un schéma-bloc ?
\item comment détermine-t-on la sortie d\textquotesingle un modèle à n entrées~?
\item un système asservi est-il toujours bouclé ? La réciproque est-elle vraie ? Donner un exemple.
\item donner les fonctions de transfert des modèles proportionnel, dérivateur et intégrateur puis donner les graphes représentant leur réponse à un échelon ;
\item donner la fonction de transfert d\textquotesingle un modèle du 1er ordre ainsi que ses paramètres caractéristiques, puis donner le graphe représentant sa réponse à un échelon. Indiquer les points caractéristiques sur ce graphe (temps de réponse, variation totale de la sortie) ;
\item donner la fonction de transfert d\textquotesingle un modèle du 2ème ordre ainsi que ses paramètres caractéristiques, puis donner le graphe représentant sa réponse temporelle à un échelon. Indiquer les points caractéristiques sur ce graphe (variation totale de la sortie...) ;
\item dans quel cas la sortie d'un modèle du 2ème ordre présente-t-elle des oscillations amorties~?
\item combien y a-t-il de dépassement \textgreater1\% pour le cas z=0,69 ? Donner leur valeur ?
\item donner l'expression qui permet de quantifier un dépassement relatif~;
\item donner les expressions de la période et de la pseudo-pulsation pour un 2ème ordre dont la réponse est oscillatoire amortie~;
\item à quels instants, les 3 premiers dépassements s'effectuent pour un 2ème ordre dont la réponse est oscillatoire amortie ?
\item on désire un temps de réponse le plus faible possible pour un modèle du 2ème ordre. Que faut-il faire~? 2 cas sont à envisager~;
\item comment détermine-t-on le temps de réponse pour un modèle du 2ème ordre~?
\item que vaut le temps de réponse réduit pour z=0,69 et z=1~? Quel est son unité ?
\item donner la fonction de transfert d'un modèle à retard.
\item expliquer la méthode d\textquotesingle identification d\textquotesingle un modèle de comportement du 1er ordre ;
\item expliquer la méthode d\textquotesingle identification d\textquotesingle un modèle de comportement du 2ème ordre oscillatoire ;
\item expliquer la méthode d\textquotesingle identification d\textquotesingle un modèle de comportement du 2ème ordre apériodique.
\item donner les caractéristiques d\textquotesingle un comportement fréquentiel ;
\item donner les méthodes analytique et graphique permettant d'obtenir la(les) pulsation(s) de coupure à -3dB.
\item comment peut-on évaluer les comportements en stabilité, rapidité et précision par lecture d'un diagramme de Bode~?
\item donner les diagrammes de Bode des fonctions de transfert des systèmes proportionnel, intégrateur, intégrateur double, dérivateur, dérivateur double, 1er ordre et 2ème ordre. Indiquer les points caractéristiques sur ces graphes.
\item expliquer la méthode d'identification d'un modèle de comportement à partir d'un diagramme de Bode.
\end{enumerate}
